\documentclass[10pt,a4paper]{amsart}
\usepackage[margin=19mm]{geometry}
\usepackage{amsmath,amssymb,mathtools}
\usepackage[scr=rsfs,cal=euler]{mathalfa}
\usepackage{xfrac}
\usepackage[unicode,hidelinks,bookmarks=false]{hyperref}
\newcommand{\ie}{i.e.\ }
\newcommand{\cf}{cf.\ }
\newcommand{\resp}{respectively\ }
\newcommand{\Subsection}{\subsection}
\providecommand{\nobreakdash}{\nobreak\hbox{-}}
\setlength{\emergencystretch}{2em}
\multlinegap0pt
\usepackage{amssymb}
\usepackage[all]{xy}
\usepackage{mathtools}
\usepackage{xcolor}
\newtheorem{theorem}{Theorem}
\newtheorem{coro}{Corollary}
\newcommand{\C}{\mathbb{C}}
\newcommand{\R}{\mathbb{R}}
\DeclareMathOperator{\U}{U}
\newcommand{\bA}{\mathbb{A}}
\newcommand{\cA}{\mathcal{A}}
\newcommand{\cH}{\mathcal{H}}
\newcommand{\cP}{\mathcal{P}}
\newcommand{\cU}{\mathcal{U}}
\DeclareMathOperator{\cusp}{cusp}
\newcommand{\valP}[1]{\left|#1\right|}
\title{The global Gan--Gross--Prasad conjecture for Fourier--Jacobi periods on unitary groups \\ III: Proof of the main theorems}
\author{Paul Boisseau, Weixiao Lu, Hang Xue}
\date{Source: arXiv:2601.01738v1 (CC BY 4.0)}
\begin{document}
\maketitle

\noindent\emph{Source and licence.} The global Gan--Gross--Prasad conjecture for Fourier--Jacobi periods on unitary groups \\ III: Proof of the main theorems. Version arXiv:2601.01738v1 (CC BY 4.0), distributed by arXiv under the Creative Commons Attribution 4.0 International licence, http://creativecommons.org/licenses/by/4.0/, as recorded in the arXiv licence metadata for this exact version and shown on the abstract page https://arxiv.org/abs/2601.01738v1. Full licence notice: \texttt{licenses/arxiv-2601.01738-CC-BY-4.0.txt}.\par\smallskip

\noindent\textit{Editorial scope.} Counted unit: the article's theorem thm:GGP\_arbitrary\_corank (source0.tex lines 3214--3223). The article's complete original proof of it is delivered with it (source0.tex lines 3225--3250, one \texttt{proof} environment of 26 lines). No mathematics is written, repaired or re-derived: the statement and the proof are the article's own lines, in the article's own notation, and no retained line is substituted or re-flowed.  Retained uncounted as the source-local context the proof needs: lines 2509--2516; lines 3183--3190.  Nothing was omitted from any retained span, so the package carries no omission record.  No retained line contains a citation, so the wrapper carries no bibliography and no bibliography entry.  No retained line contains a figure environment, an image-inclusion command or drawing code, so no image is delivered and the package declares no asset dependency.  Editorial formatting only, in addition: an AMS \texttt{amsart} preamble that restates the article's own declarations and macros used by the retained lines, namely the environments \texttt{theorem}, \texttt{coro} on separate counters, together with the standard packages those lines need.  The retained environments print in this document as Theorem 1, Coro 1 in order of appearance, whereas the article prints the same items with the numbering of its own full document.  The complete native source of the article as distributed in the arXiv e-print for this exact version is bundled unchanged as \texttt{sources/192105/source0.tex}.
\begin{theorem} 
    \label{thm:ggp_intro}
    The following are equivalent:
    \begin{enumerate}
    \item $L(\frac{1}{2}, \Pi_\lambda \otimes  \overline{\mu}) \not=0$; 
    \item there exist $V \in \cH$, a standard parabolic subgroup $Q \subset \U_V$ and $\sigma \in \Pi_{\cusp}(M_Q)$ such that $\Pi$ is the weak base change of $(Q,\sigma)$ and the Fourier--Jacobi period $\cP_{\U_V'}(\cdot,\lambda)$ does not vanish identically on $\cA_{Q,\sigma}(\U_V) \otimes \omega^\vee$.
    \end{enumerate}
    \end{theorem}

\begin{coro}
\label{cor:non_vanishing_of_periods}
    The following are equivalent.
    \begin{enumerate}
        \item There exist $\varphi_V \in \sigma_V$, $\varphi_W \in \sigma_W$  and $\phi \in \nu^\vee$ such that $\cP_{\cH_W} \left( \varphi_V \otimes \varphi_W , \phi \right) \neq 0$.
        \item There exist $\varphi_V \in \sigma_V$, $\varphi_\Sigma \in \Sigma$, $\Phi \in \omega^\vee_{V}$ and $s \in \C$ such that $E_{P(X)}^{\U(V)}(\varphi,.)$ has no pole at $s$ and $\cP_{\U'_V}\left( \varphi_V \otimes E_{P(X)}^{\U(V)}(\varphi_\Sigma,s) , \Phi \right) \neq 0$.
    \end{enumerate}
\end{coro}

\begin{theorem} \label{thm:GGP_arbitrary_corank}
Let $\Pi$ be a discrete Hermitian Arthur parameter of $G_n \times G_m$. The
following assertions are equivalent:
    \begin{enumerate}
        \item $L(\frac{1}{2}, \Pi \otimes \overline{\mu}) \neq 0$;
        \item there exist non-degenerate skew-Hermitian spaces $W\subset V$ and $\sigma$ a cuspidal automorphic representation of $\U(V)(\bA) \times \U(W)(\bA)$ such that $W^\perp$ is split, the weak base change of $\sigma$ to
            $G_n \times G_m$ is $\Pi$ and $\cP_{\cH_W}$ does not vanish identically on $\sigma \otimes
            \nu^\vee$.
    \end{enumerate}
\end{theorem}

\begin{proof}
    Let $\alpha_1, \hdots, \alpha_r$ be automorphic characters of $\bA_E^\times$ such that the $\alpha_1, \hdots, \alpha_r, \alpha_1^*, \hdots, \alpha_r^*$ are two by two distinct. Let $Q_n$ be the standard parabolic of $G_n$ with Levi factor $G_1^{2r} \times G_m$, and define for all $s \in i \R$
    \begin{equation*}
        \widetilde{\Pi}_s:= \Pi_n \boxtimes I_{Q_n}^{G_n}( \alpha_1\valP{.}^s_E \boxtimes \alpha_1^*\valP{.}^{-s}_E \boxtimes \hdots \boxtimes \alpha_r\valP{.}^s_E \boxtimes \alpha_r^*\valP{.}^{-s}_E \boxtimes \Pi_m ).
    \end{equation*}
    Note that the $\widetilde{\Pi}_s$ are $(G,H,\overline{\mu})$-regular Hermitian Arthur parameters. By elementary properties of the Rankin--Selberg $L$-function, assertion (1) of Theorem~\ref{thm:GGP_arbitrary_corank} is equivalent to the following assertion.
    \begin{enumerate}
        \item[(1')] There exists an $s \in i\R$ such that $L(\frac{1}{2},\widetilde{\Pi}_s) \neq 0$. 
    \end{enumerate}
    By the Gan--Gross--Prasad conjecture for $(G,H,\overline{\mu})$-regular Hermitian Arthur parameters from Theorem~\ref{thm:ggp_intro}, this implies (2). 
    
    Conversely, let $\sigma=\sigma_V \boxtimes \sigma_W$ be a cuspidal automorphic representation of $\cU_W$ whose weak base change to $G_n \times G_m$ is $\Pi$. Let $P$ be the parabolic subgroup of $\U(V)$ stabilizing the flag $E x_1 \subset \hdots \subset X$, with Levi factor $T_r \times \U(W)$. 
    \begin{equation}
        \label{eq:kappa_defi}
    \kappa:=\alpha_1 \boxtimes \hdots \boxtimes \alpha_r,
    \end{equation}
    which is an automorphic character of $T_r(\bA)$.
    Then for all $s \in i \R$ the representation $\widetilde{\Pi}_s$ is the weak base change of $(\U(V) \times P,  \sigma_V \boxtimes \kappa_s \boxtimes \sigma_W)$.
    
    Define $\tau=I_{B_r}^{G_r}(\kappa)$ so that $I_P^{\U(V)}(\kappa_s \boxtimes \sigma_W)=I_{P(X)}^{\U(V)} (\tau_s \boxtimes \sigma_W)$. Then by Corollary~\ref{cor:non_vanishing_of_periods}, (2) of Theorem~\ref{thm:GGP_arbitrary_corank} is equivalent to the following assertion.
    \begin{enumerate}
        \item[(2')] There exists an $s \in i \R$ such that the bilinear form $\cP_{\U'_V}$ is non-zero on $\sigma_V \otimes (I_{P(X)}^{\U(V)} \tau_s \boxtimes \sigma_W) \otimes \omega^\vee_{V}$.
    \end{enumerate}
    But (2') implies (1') by Theorem~\ref{thm:ggp_intro}. This concludes the proof of Theorem~\ref{thm:GGP_arbitrary_corank}.

\end{proof}

\end{document}
